\SourceLeaf{643}
\section*{BEMERKUNGEN.}
\addcontentsline{toc}{section}{Bemerkungen}
\par\nopagebreak[4]\smallskip\noindent\makebox[\linewidth][c]{\rule{.08\linewidth}{.35pt}}\par

Die hier unter 21 Nummern zusammengestellten, bruchstückhaft aufgezeichneten Untersuchungen stammen aus recht weit auseinanderliegenden Zeiten. Nr. 1 und 2 finden sich in einem Tagebuch zwischen den Protokollen von Beobachtungen, die im März, Juni und Juli 1833 zu den durch Magnete induzierten galvanischen Strömen angestellt wurden. Nr. 6, 12, 18 und 21 stehen auf gesonderten Blättern. Außer Nr. 18, die ein Datum trägt, lassen sie sich nicht genauer datieren. Die übrigen Nummern, mit Ausnahme des Briefes [20.], sind hier in derselben Reihenfolge wiedergegeben wie in einem Notizbuch, in dem allerdings zahlreiche Ausführungen zu ganz anderen Themen zwischen ihnen stehen. Die letzte dieser Nummern mit dem Beweis des Fundamentalsatzes von \textsc{Ampère} wurde erst nach 1843 eingetragen. Die anderen scheinen aus der Zeit von 1833 bis 1836 zu stammen.

Die verschiedenen Formen, die hier für das Gesetz der Wechselwirkungen zwischen galvanischen Stromelementen angenommen werden, folgen alle aus dem vor allem in Nr. 20 hervorgehobenen Prinzip, dass Magnetismus und galvanische Ströme austauschbar sind. Die in diesem Brief angesprochenen Untersuchungen von \textsc{Wilhelm Weber} bilden die Vorarbeiten zu einer Theorie, deren Aufstellung 1846 in der ersten Abhandlung über elektrodynamische Maßbestimmungen vollendet wurde. Nach dieser Theorie wird die gesamte Wechselwirkung zwischen zwei elektrischen Teilchen (wobei $e$, $e'$ mit den entsprechenden Vorzeichen versehen sind) im gegenseitigen Abstand $r$ durch
\[
ee'\left(\frac1{rr}-\frac1{cc}\frac1{rr}\frac{\mathrm dr^2}{\mathrm dt^2}+\frac2{cc}\frac1r\frac{\mathrm ddr}{\mathrm dt^2}\right)
\]
gemessen. Ein positiver Wert dieser Größe bedeutet Abstoßung, ein negativer Anziehung. In dem Ausdruck bezeichnet $t$ die Zeit und $c$ eine Geschwindigkeit. Bei ihren Untersuchungen (1855) zur Rückführung der Stromstärkemessungen auf mechanisches Maß fanden \textsc{Kohlrausch} und \textsc{Weber} für diese Geschwindigkeit den Wert $439450.10^6\,\dfrac{\text{Millimeter}}{\text{Sekunde}}$.

\SourceLeaf{644}

Der Lehrsatz in Nr. 4 hat eine rein geometrische Form erhalten. Wegen seiner Bedeutung für die Theorie der galvanischen Ströme glaubte ich, ihm hier seinen Platz geben zu müssen. Das Integral, das die Anzahl der Umschlingungen der geschlossenen Kurve $s$ mit dem System geschlossener Kurven $s'$ bestimmt, ergibt, wenn man für $\mathrm ds'$ die Elemente aller im Raum vorhandenen geschlossenen galvanischen Ströme einsetzt, die algebraische Summe der Stärken derjenigen Ströme, die eine von $s$ begrenzte, im Übrigen aber beliebig gewählte Fläche durchdringen. Bei dieser Deutung der $\mathrm ds'$ ist das Integral selbst jedoch gleich $\dfrac1{4\pi}\int\dfrac{\mathrm dV}{\mathrm ds}\mathrm ds$, wenn $V$ wie in Nr. 9 die Potenzialfunktion für die magnetischen Wirkungen der Ströme $s'$ bezeichnet. Der Satz entspricht also dem Satz, den \textsc{Gauss} in seiner Abhandlung über die Anziehung der Ellipsoide aufgestellt hat. Dieser bestimmt die Masse innerhalb einer geschlossenen Fläche $(\omega)$ aus den nach innen gerichteten Normalkräften ihrer Anziehung $\left(\dfrac{\mathrm dV}{\mathrm dN}\right)$ an der Fläche durch $\dfrac1{4\pi}\int\dfrac{\mathrm dV}{\mathrm dN}\mathrm d\omega$.

Den angegebenen Wert des obigen Integrals $\dfrac1{4\pi}\allowbreak\int\dfrac{\mathrm dV}{\mathrm ds}\mathrm ds$ erhält man zum Beispiel, indem man die Integralausdrücke für die Ableitungen von $V$ nach den Koordinaten in Integrale über Flächen umwandelt, die von den einzelnen Stromleitern $s'$ begrenzt werden, wobei die Flächen $\omega$ beliebig gewählt sind. Ein in Art. 38 der allgemeinen Theorie des Erdmagnetismus angedeuteter Satz gibt den Unterschied der Werte der Potenzialfunktion an entsprechenden Stellen auf beiden Seiten einer Fläche an, die auf beiden Seiten in geeigneter Weise mit entgegengesetztem magnetischem Fluidum belegt ist. Nach diesem Satz zeigt die so erhaltene Form der Ableitung von $V$ nach $s$ unmittelbar, dass das gesuchte Integral gleich der algebraischen Summe der Stärken der Ströme ist, die in den Begrenzungslinien der von der Kurve $s$ durchsetzten Flächen $\omega'$ fließen.

Integrale über eine geschlossene Kurve $s$ lassen sich mit dem folgenden Satz in Integrale über eine von $s$ begrenzte Fläche $\omega$ umwandeln. Für beliebige rechtwinklige gerad- oder krummlinige Koordinaten $\xi,\eta,\zeta$ wird das Quadrat des Längenelements allgemein durch einen Ausdruck der Form
\[
\xi'\xi'\mathrm d\xi^2+\eta'\eta'\mathrm d\eta^2+\zeta'\zeta'\mathrm d\zeta^2
\]
dargestellt. Sind an den Punkten der Fläche $\omega$ beliebige Funktionen $\lambda,\mu,\nu$ der Koordinaten $\xi,\eta,\zeta$ mitsamt ihren ersten Ableitungen stetig veränderlich, so gilt stets die Gleichung
\[
\begin{aligned}
&\int\left(\lambda\frac{\mathrm d\xi}{\mathrm ds}+\mu\frac{\mathrm d\eta}{\mathrm ds}+\nu\frac{\mathrm d\zeta}{\mathrm ds}\right)\mathrm ds\\
&=\int\left\{\left(\frac{\mathrm d\mu}{\mathrm d\zeta}-\frac{\mathrm d\nu}{\mathrm d\eta}\right)\frac{\xi'}{\eta'\zeta'}\frac{\mathrm d\xi}{\mathrm dn}+\left(\frac{\mathrm d\nu}{\mathrm d\xi}-\frac{\mathrm d\lambda}{\mathrm d\zeta}\right)\frac{\eta'}{\zeta'\xi'}\frac{\mathrm d\eta}{\mathrm dn}\right.\\
&\qquad\left.+\left(\frac{\mathrm d\lambda}{\mathrm d\eta}-\frac{\mathrm d\mu}{\mathrm d\xi}\right)\frac{\zeta'}{\xi'\eta'}\frac{\mathrm d\zeta}{\mathrm dn}\right\}\mathrm d\omega
\end{aligned}
\]
unter der Voraussetzung, dass $\xi,\eta,\zeta$ im ersten Integral die Koordinaten eines Punkts des Längenelements $\mathrm ds$, im zweiten $\xi,\eta,\zeta$ die Koordinaten eines Punkts des Flächenelements $\mathrm d\omega$ und $n$ die Normale zu diesem Flächenelement bezeichnen. Die positive Richtung der Normalen ist folgendermaßen zu wählen. Es sei $\mathrm dt$ das erste Element einer Kurve, die von einem Punkt des Elements $\mathrm ds$ ausgeht, zu diesem Element senkrecht steht und in der Fläche $\omega$ liegt. Dann müssen die positiven Richtungen von $\xi'\mathrm d\xi$, $\eta'\mathrm d\eta$, $\zeta'\mathrm d\zeta$ durch stetige Verschiebung des Koordinatensystems im Raum der Reihe nach mit den positiven Richtungen von $\mathrm dn$, $\mathrm ds$, $\mathrm dt$ zur Deckung gebracht werden können.

\SourceLeaf{645}

Die wiederholte Anwendung dieses Satzes liefert auch den Beweis des Fundamentalsatzes von \textsc{Ampère} in der allgemeinen Form, in der die Potenzialfunktion für die Wechselwirkung zwischen den auf bestimmte Weise mit magnetischem Fluidum belegten Flächen unmittelbar auf die Potenzialfunktion für die Wechselwirkung zwischen galvanischen Strömen zurückgeführt wird, deren Lage und Stärke durch diese Flächen und die Magnetisierung bestimmt sind.

Dem Beweis in Nr. 9 für die Gleichheit der Werte der verschiedenen Ausdrücke für die Potenzialfunktion $V$ kann man eine symmetrische Form geben, indem man die Funktion $R$
\[
\begin{aligned}
={}&xx\operatorname{arctang}\frac{yz}{xr}+yy\operatorname{arctang}\frac{zx}{yr}+zz\operatorname{arctang}\frac{xy}{zx}\\
&+2yzi\operatorname{arctang}\frac{xi}r+2zxi\operatorname{arctang}\frac{yi}r+2xyi\operatorname{arctang}\frac{zi}r
\end{aligned}
\]
einführt, in der $i$ für $\BareRoot{-1}$ steht, und berücksichtigt, dass die Gleichungen
\[
\begin{array}{ll}
\dfrac{\mathrm ddR}{\mathrm dx^2}=2\operatorname{arctang}\dfrac{yz}{xr},&\dfrac{\mathrm ddR}{\mathrm dy\mathrm dz}=2i\operatorname{arctang}\dfrac{xi}r\\[5pt]
\dfrac{\mathrm ddR}{\mathrm dy^2}=2\operatorname{arctang}\dfrac{zx}{yr},&\dfrac{\mathrm ddR}{\mathrm dz\mathrm dx}=2i\operatorname{arctang}\dfrac{yi}r\\[5pt]
\dfrac{\mathrm ddR}{\mathrm dz^2}=2\operatorname{arctang}\dfrac{xy}{zr}.&\dfrac{\mathrm ddR}{\mathrm dx\mathrm dy}=2i\operatorname{arctang}\dfrac{zi}r
\end{array}
\]
\[
\frac{\mathrm ddR}{\mathrm dx^2}+\frac{\mathrm ddR}{\mathrm dy^2}+\frac{\mathrm ddR}{\mathrm dz^2}=(2m+1)\pi,\qquad\frac{\mathrm d^3R}{\mathrm dx\mathrm dy\mathrm dz}=-2\frac1r
\]
gelten. Mit den Ableitungen der Funktion $R(x,y,z)$ lassen sich auch die Potenzialfunktionen für die magnetische Wirkung galvanischer Ströme in endlicher Form darstellen, deren Leiter aus geradlinigen, zu den Achsen eines rechtwinkligen Koordinatensystems parallelen Abschnitten bestehen.

Die Kenntnis der Anordnung der Längenvergleiche, mit der \textsc{Gauss} einen Abschnitt eines geteilten Maßstabs in Teilen des gesamten Maßstabs bestimmte, verdanken wir der Aufzeichnung, die sich Herr Geheimer Hofrat \textsc{Weber} im Jahr 1839 oder 1840 gemacht hat.

Die theoretischen Untersuchungen zu den Beugungserscheinungen wurden wahrscheinlich durch das Werk zu diesem Thema angeregt, das \textsc{F. M. Schwerd} 1835 herausgab. Die beiden allgemeinen Formeln für die Wirkung eines leuchtenden Punkts $P$ auf einen Punkt $p$ sind verschieden. Man betrachte Flächenintegrale, deren Elemente von der Lage der Ebene unabhängig sind, die durch einen Punkt des zugehörigen Flächenelements $\mathrm ds$ und durch die Punkte $P$ und $p$ geht. Bei den zugehörigen Kurvenintegralen sind die Elemente ebenfalls von der Lage der Ebene unabhängig, die durch einen Punkt des zugehörigen Elements der Begrenzungslinie $u$ und durch die Punkte $P$ und $p$ geht. Ihre Differenziale bezeichnen aber nur eine Änderung des Winkels $\theta$, den jene Ebene mit einer festen Ebene durch $P$ und $p$ einschließt. Die allgemeine Umwandlung dieser Flächenintegrale in die Kurvenintegrale ergibt sich aus der Gleichung
\SourceLeaf{646}
\[
\begin{aligned}
\int Q\mathrm d\theta=\int\frac{h.Q.\sin v.\mathrm du}{rR\sin w}
={}&\int\frac h{rR}.\frac1{2\sin\frac12w^2}.\frac{\mathrm dQ}{\mathrm d(R-r)}\frac{\mathrm d(R+r)}{\mathrm d\rho}.\mathrm ds\\
&-\int\frac h{rR}.\frac1{2\cos\frac12w^2}.\frac{\mathrm dQ}{\mathrm d(R+r)}.\frac{\mathrm d(R-r)}{\mathrm d\rho}\,\mathrm ds
\end{aligned}
\]
Diese Gleichung ist ein Spezialfall des in der vorangehenden Bemerkung genannten Satzes, wenn $Q$ eine nur von $R$ und $r$ abhängige Größe bezeichnet, die zusammen mit ihren partiellen Ableitungen nach $R+r$ und $R-r$ an den Punkten der Fläche $s$ stetig veränderlich ist. Außerdem ist $w$ genauer als im Text bestimmt. Es bezeichnet den Winkel zwischen $R$ und $r$, den die positiven Richtungen dieser Linien einschließen, wobei diese Richtungen mit denen des fortschreitenden Lichtstrahls übereinstimmend gewählt sind.

\begin{flushright}\textsc{Schering.}\end{flushright}

\par\nopagebreak[4]\medskip\noindent\makebox[\linewidth][c]{\rule{.20\linewidth}{.35pt}}\par
\clearpage
